\subsection{Passenger journeys, vehicle legs, and observed events}
\label{sec:journey-leg-event-contract}

The OD demand $f_q^t$ is interpreted as \emph{passenger journey demand}.  A
journey begins when a passenger first boards the public transport system at
$s_o$ and ends when that passenger finally alights at $s_d$.  It may contain
one or more \emph{vehicle legs}.  A vehicle leg is the movement on one vehicle
trip between a boarding event and its corresponding alighting event.  Two
consecutive legs are connected by a transfer; a journey with $L$ legs has
$L-1$ transfers.  Boarding the first vehicle is not a transfer.

This distinction is essential when boarding and alighting counts are used for
OD estimation.  One unit of demand assigned to an $L$-leg journey contributes
one boarding and one alighting on every leg.  It therefore generates one
initial boarding, one final alighting, $L-1$ transfer alightings, and $L-1$
transfer boardings.  A transfer boarding is not a new journey origin, and a
transfer alighting is not a final journey destination.  Consequently, total
boardings and alightings are generally not journey-origin and
journey-destination marginals in a network with transfers.

The atomic boarding or alighting observation is associated with one timetable
event identified by measurement method, event type, stop, clock time, and a
trip or an unambiguous line reference.  Platform, physical-stop, route-period,
or other aggregates are derived views and require an explicit, versioned
mapping.  The initial reduced-dimensional backend accepts boarding and
alighting events; the existing \texttt{load} schema value is outside that
backend until a corresponding response contract is implemented and tested.

An absent record denotes an unobserved event and creates no likelihood term.  A
present record whose value is zero is an observed zero and remains in the
likelihood.  Duplicate, negative, non-finite, unmatched, and ambiguous records
are errors; APC cleaning, outage exclusion, and imputation occur before the
immutable observation table is constructed and are recorded as input
provenance.

The external OD key remains the triplet comprising
\path{origin_stop_id}, \path{dest_stop_id}, and \path{time_bin_id}.  A
platform-to-physical-stop mapping may be used for transfer construction, but it
is an explicit, fingerprinted input.  An authoritative operator or GTFS
parent-station mapping is preferred; a generated mapping requires review and
cannot be accepted silently.  Published output remains at scenario-stop level
unless an explicit aggregation policy requests physical stops.

For a reduced conditional destination model, let $Q_{ot}$ denote the number of
journeys whose first boarding occurs at origin $o$ in desired-departure bin
$t$.  It must not be set equal to all boarding counts at $(o,t)$ when transfer
boardings may be present.  The first reduced implementation therefore admits
either externally provided journey-entry totals or productions estimated
through a declared low-dimensional nonnegative basis with regularization.  It
does not provide a raw-boarding production mode.  Which productions,
attractiveness terms, route shares, detection parameters, and dispersion
parameters are fixed, normalized, estimated, or regularized must be declared
by every model specification.

These definitions constitute the observation contract for the planned
reduced-dimensional backend.  They do not redirect or replace the existing
detailed assignment and estimation APIs.  The time-expanded assignment remains
the explicit validation backend and is not evaluated inside the reduced
objective or gradient.

